% 総目録（紙見本・解説記事編）．surikumi-kaisetsu の部品を，
% 実際の解説記事と同じ版面で組んだ姿と命令名を並べて示す．
% surikumi 本体の部品は版面が異なるため surikumi-catalog.pdf にある．
% 本文の例は本見本のために書き下ろしたものであり，実在の記事の引用ではない．
\documentclass[lualatex,paper=b5j,fontsize=10pt,twoside]{jlreq}

\usepackage{surikumi-kaisetsu}

% --- 見本用の局所部品 -------------------------------------------------
% カタログとしての注釈であり，紙面部品そのものではない．記事本体では使わない．
\newcommand{\CatName}[1]{{\ttfamily\fontsize{7.6}{9}\selectfont #1}}

\newcommand{\CatLabel}[2]{%
  \par\addvspace{.5\baselineskip}%
  \noindent{\MKHeadingFace\fontsize{7.6}{9}\selectfont
    \color{SMrule}#1}\hspace{.6em}\CatName{#2}\par
  \nobreak\vspace{.5mm}%
  \noindent{\color{SMpale}\rule{\linewidth}{.4pt}}\par
  \nobreak\addvspace{.25\baselineskip}%
}

\newcommand{\CatUse}[1]{%
  \par\addvspace{.15\baselineskip}%
  \noindent{\color{SMrule}\fontsize{7.6}{10}\selectfont #1}\par
  \nobreak\addvspace{.2\baselineskip}%
}

% 扉の飾り．参考紙面の図案を写さず，TeX で生成した独自の図形を渡す．
\newcommand{\CatOrnament}{%
  \begin{tikzpicture}[x=1mm,y=1mm]
    \foreach \i in {0,...,6}{
      \draw[SMink,line width=.35pt]
        (0,\i*4) -- (30,\i*4);
    }
    \foreach \i in {0,...,6}{
      \draw[SMink,line width=.9pt]
        (\i*5,0) .. controls (\i*5+2.5,12) .. (\i*5,24);
    }
  \end{tikzpicture}%
}

\MathKaisetsuSetup{
  journal-name={数理組版},
  masthead-first={SURIKUMI},
  masthead-second={REVIEW},
  issue-date={August 2026},
  issue-number={Number 8},
  masthead=true,
  align=center,
  section-style=bar,
  feature={総目録／解説記事の部品},
  title={解説記事の部品見本},
  subtitle={},
  author={\MKName{数理}{太郎}},
  author-reading={すうり・たろう},
  affiliation={数理組版編集部}
}

\begin{document}

\MKMakeTitle
\MKArticleReset

この見本は \CatName{surikumi-kaisetsu.sty} が用意する部品を，
実際の記事と同じ版面――本文 9\,pt，行送り 14.2\,pt，1 段 45 行――で
一通り並べたものである．
上の扉そのものが \CatName{\string\MKMakeTitle} の \CatName{align=center} の例であり，
号の巻頭記事に使う型である．

灰色の小見出しと \CatName{等幅の命令名} は，この見本のための注釈であって，
紙面部品ではない．

命令には短縮名がある．\CatName{\string\MathKaisetsuSetup} は
\CatName{\string\MKSetup}，\CatName{\string\MKMakeTitle} は
\CatName{\string\MKTitle}，\CatName{\string\MKArticleReset} は
\CatName{\string\MKReset} である．どちらで書いても同じものが組まれる．
全対応表は \CatName{docs/catalog.md} の第13章にある．

\MKSection{見出しの四つの型}

節見出しには四つの型がある．紙面全体の既定は
\CatName{section-style} で決め，例外だけを
\CatName{\string\MKSection[型]} で指定する．
段落ごとの気分で型を変えない．
いま読んでいるこの見出しが既定の \CatName{bar} である．

\MKSubsection{小見出し}

小見出しは 1 字下げの位置から始まり，番号は節番号と組にして
\CatName{1.1} の形で付ける．直後の段落も 1 字下げで始まる．
番号を付けない見出しには星付き形式を使う．

\MKSection[pillars]{pillars ――特集記事の節見出し}

両端の柱の間に中央寄せで置く型である．
特集の中の記事で使う．

\MKSection[band]{band ――章が切り替わる位置}

上下の罫と両端の柱を持つ，最も強い型である．
記事の中で話題が大きく変わる位置に限って使う．

\MKSection[plain]{plain ――罫も柱もなし}

短い記事や囲みの中で使う．

\MKSection{行頭の見出しと注意}

\CatLabel{行頭の見出し}{\string\MKRunIn}
\MKRunIn{注意}行頭にゴシック体の語を置き，本文をそのまま続ける．
節を立てるほどではない切り替えに使う．

\CatLabel{注意などの段落}{\string\MKRemark，mkremark}
\MKRemark{「最小」という言い方は歴史的なものである．
停留値であれば足りる場面が多い．}

\begin{mkremark}[定義]
複数の段落や別行数式を含むときは環境の形を使う．
省略可能な引数で行頭の語を変えられる．
\end{mkremark}

\CatLabel{引用}{mkquote}
\begin{mkquote}
引用は左を 1 字下げる．出典は引用のあとの本文で示す．
\end{mkquote}

\MKSection{数式番号}

既定では節ごとに番号を付け，節が変わると
\CatName{(1.1)} から数え直す．いまは第 4 節なので，次の式は
\CatName{(4.1)} になる．

\begin{equation}
  S[x]=\int_{t_a}^{t_b}L\bigl(\dot{x}(t),x(t)\bigr)\,\mathrm{d}t
  \label{eq:action}
\end{equation}

番号を付けるのは本文中で参照する式だけである．
式\eqref{eq:action} のように参照しない式は
\CatName{\string\[...\string\]} で組む．

\[
  \dfrac{\partial L}{\partial x}
  -\dfrac{\mathrm{d}}{\mathrm{d}t}\dfrac{\partial L}{\partial\dot{x}}=0
\]

\CatUse{通し番号にするときは \string\MKEquationNumbering\{continuous\}，
番号を付けないときは \{none\} とする．}

Σ の扱いは \CatName{surikumi} 本体と同じで，行内式
$S_n=\sum_{k=1}^{n}a_k$ と別行数式

\[
  S_n=\sum_{k=1}^{n}a_k
\]

で本体の大きさが等しく，上下限はどちらも真上と真下に置かれる．

\MKSection{注と参考文献}

段の下に置く注は \CatName{*1)} の形で番号を付ける\footnote{これが注である．
折り返しは番号の右にそろう．}．
本文中の参考文献番号 \CatName{1)} とは形が違うため，
同じ段に並んでも取り違えない\footnote{二つ目の注．番号は記事ごとに
\CatName{\string\MKArticleReset} で 1 から数え直す．}．

本文に印を出さない注もある．

\MKNote{本稿の図はすべて TeX で生成した．}

本文からは肩付きの番号で参照する．
詳しくは文献\MKCite{\ref{ref:variational}} を見る．
複数を並べるときは \MKCite{1,2} のようになる．

\MKSection{図と表}

図の見出しは図の下，表の見出しは表の上に置く．
番号はゴシック体，本文は明朝体で，段の幅いっぱいに組む．

\begin{figure}[t]
  \begin{mkfig}
    \begin{tikzpicture}[x=1mm,y=1mm]
      \draw[SMink,line width=.5pt] (0,0) -- (46,0);
      \draw[SMink,line width=.9pt]
        (0,2) .. controls (16,20) and (30,20) .. (46,2);
      \draw[SMrule,line width=.4pt,dashed]
        (0,2) .. controls (16,12) and (30,26) .. (46,2);
      \fill[SMink] (0,2) circle (.7);
      \fill[SMink] (46,2) circle (.7);
    \end{tikzpicture}
  \end{mkfig}
  \caption{実線が停留経路，破線が近傍の経路である．
    図の中身は \CatName{mkfigurebody ／ mkfig} で段の中央にそろえる．}
  \label{fig:paths}
\end{figure}

図\ref{fig:paths} のように，本文からは通常どおり参照する．

\MKSection*{番号を付けない節}

星付き形式は番号を進めない．付録や補足に使う．

\begin{mkrefs}
  \item 変分原理に関する解説．\label{ref:variational}
  \item 解析力学の教科書．
\end{mkrefs}

\MKByline

% =====================================================================
% 2 本目の記事．扉の型を align=left に変え，副題と飾りを付ける．
% \MKMakeTitle が前の記事を締めるので，直前に \clearpage は置かない．
\MathKaisetsuSetup{
  masthead=false,
  align=left,
  feature-bar=true,
  section-style=pillars,
  feature={総目録／扉の型},
  title={特集記事の扉},
  subtitle={左寄せの表題と副題，外側の飾り},
  author={\MKName{数理}{二郎}},
  author-reading={すうり・じろう},
  affiliation={数理組版編集部},
  ornament={\CatOrnament},
  ornament-width={34mm}
}

\MKMakeTitle
\MKArticleReset

この扉が \CatName{align=left} の例である．
細罫と太罫で挟んだ特集名，左寄せの表題，副題，署名を置き，
外側に飾りを入れられる．
特集の中に置く記事に使う．

飾りは文書側で用意する．参考紙面の図案を写さず，
\CatName{\string\includegraphics} か，TeX で生成した独自の図形を
\CatName{ornament} に渡す．この扉の飾りも本見本のために描いたものである．

% 段抜きの囲みは浮動体である．紙面の下に置くには，その紙面の組版が
% 終わる前に宣言しておく必要があるため，ここで declare しておく．
\begin{mkfull}[b]
\MKRunIn{段抜きの囲み}\CatName{mkfullinsert} は両段にまたがる囲みで，
既定では紙面の下に置く．省略可能な引数で位置を指定できる．
節を立てずに独立した話題を差し込むときに使う．
浮動体なので，置きたい紙面の組版が終わる前に書いておく．
\end{mkfull}

\MKSection{節見出しは pillars にしてある}

この記事では \CatName{section-style=pillars} を既定にした．
記事ごとに \CatName{\string\MathKaisetsuSetup} で組み直せる．

\MKSection{年表と囲み}

\CatLabel{年表}{mkchronology ／ mkchrono，\string\MKChronoItem}
\begin{mkchrono}{変分法の歩み}
  \MKChronoItem{1696}{Bernoulli が最速降下線の問題を提示する}
  \MKChronoItem{1744}{Euler が変分法の一般的な手続きをまとめる}
  \MKChronoItem{1788}{『解析力学』が刊行される}
\end{mkchrono}

\CatLabel{段の中の囲み}{mkinsert}
\begin{mkinsert}
\MKRunIn{要点}停留値であることと最小であることは別の主張である．
二階の変分を調べるまでは，どちらとも言えない．
\end{mkinsert}

\CatUse{段抜きの囲み mkfullinsert は，この紙面の下に置いてある．}

\MKSection{1 冊に複数の記事を置く}

記事ごとに扉の内容を組み直し，\CatName{\string\MKMakeTitle} と
\CatName{\string\MKArticleReset} を続けて書く．
\CatName{\string\MKMakeTitle} が前の記事を締めるので，
その直前に \CatName{\string\clearpage} は置かない．
注と数式の番号は，この記事でも 1 から数え直されている\footnote{記事が
変わって最初の注なので \CatName{*1)} に戻っている．}．

\MKByline

\end{document}
